\subsection{Step 4: Harmonization Dataset Construction}
\label{sec:step4_harmonization_dataset}

Step~4 transforms reviewed equivalence decisions into a harmonization map that can be propagated back to the full occurrence-level dataset. In the corrected workflow, this step is deliberately conservative: only pairs explicitly marked \texttt{equivalent} in Step~3 are allowed to create harmonized groups. Pairs marked \texttt{not\_equivalent} or \texttt{uncertain}, as well as pairs with blank review fields, do not create merges.

The grouping rule is connected-component closure over the reviewed positive-pair graph. Because chain merges can over-collapse the node universe when the review edges are noisy, Step~4 now exports an explicit component audit so that multi-name or multi-member groups can be checked before the harmonization layer is treated as frozen. Representative harmonized names are selected transparently by the rule ``most frequent label, then lexicographic tie-break,'' with optional manual overrides recorded separately.

If Step~3 has not yet been reviewed, Step~4 remains valid but yields a no-merge baseline. This is the current pipeline state.

\subsubsection{Outputs of this Step}
\label{sec:step4_outputs}

The corrected Step~4 produces:
\begin{itemize}
    \item \texttt{04\_textual\_instances\_harmonization\_mapping.csv}
    \item \texttt{04\_textual\_instances\_harmonized.csv}
    \item \texttt{04\_skills\_harmonized.json}
    \item \texttt{04\_skills\_harmonized.csv}
    \item \texttt{04\_unique\_skills\_harmonization\_mapping.csv}
    \item \texttt{04\_unique\_skills\_harmonized.csv}
    \item \texttt{04\_harmonization\_summary.csv}
    \item \texttt{04\_harmonized\_groups\_summary.csv}
    \item \texttt{04\_harmonized\_name\_selection.csv}
    \item \texttt{04\_component\_merge\_audit.csv}
\end{itemize}

The textual-instance mapping is the main harmonization file for downstream steps. The unique-skill summary files are retained as derived convenience outputs, but the harmonization logic itself is defined at the textual-instance level.

\subsubsection{Fields to Fill After Concluding This Step}
\label{sec:step4_fields_to_fill}

No new fields are meant to be filled inside Step~4 itself. The substantive prerequisite is still upstream:
\begin{itemize}
    \item complete the Step~3 review fields \texttt{review\_decision} and \texttt{review\_notes} in
    \texttt{03\_review\_template\_textual\_instances.csv}
\end{itemize}

After those decisions are entered, Step~4 should be rerun so that the harmonized node universe reflects reviewed equivalence rather than the current no-merge baseline.

\subsubsection{Implementation Note and Current Run}
\label{sec:step4_implementation_note}

The current corrected run reports:
\begin{itemize}
    \item \texttt{reviewed\_pairs\_equivalent = 0}
    \item \texttt{reviewed\_pairs\_not\_equivalent = 0}
    \item \texttt{reviewed\_pairs\_uncertain = 0}
    \item \texttt{review\_pairs\_without\_decision = 160}
    \item \texttt{total\_text\_instances\_before\_harmonization = 895}
    \item \texttt{total\_harmonized\_groups = 895}
    \item \texttt{multi\_text\_instance\_harmonized\_groups = 0}
    \item \texttt{text\_instances\_name\_changed = 0}
    \item \texttt{full\_skill\_records\_name\_changed = 0}
    \item \texttt{harmonization\_mode = no\_merge\_baseline}
\end{itemize}

The current Step~4 outputs are therefore methodologically valid, but they still describe the pre-review baseline. Any actual harmonization requires reviewed \texttt{equivalent} decisions to be entered in Step~3 and then propagated through a rerun of Steps~4 to~9.
